% Source: https://artofproblemsolving.com/community/c3h2357523_unpostbanned_
% Stable identity: aops-topic-2357523-q08
Look for solutions of the recurrence of the form $a_n=r^n$: then $r^n=r^{n-1}+2r^{n-2}$, that is, $r^2=r+2$, so $r=2$ or $r=-1$. Both $2^n$ and $(-1)^n$ satisfy the recurrence, and so does every combination $\alpha\cdot2^n+\beta(-1)^n$. Choose $\alpha,\beta$ to match the first two terms:
\[2\alpha-\beta=0,\qquad4\alpha+\beta=4,\]
so $\alpha=\dfrac23$ and $\beta=\dfrac43$. A sequence satisfying the recurrence is determined by its first two terms, so for all $n\ge1$
\[a_n=\dfrac23\cdot2^n+\dfrac43(-1)^n=\boxed{\dfrac{2^{n+1}+4(-1)^n}{3}}.\]
Check: the formula gives $0,4,4,12,20,44,\dots$, and indeed $a_3=4+0=4$, $a_4=4+8=12$, $a_5=12+8=20$, $a_6=20+24=44$.
